% day1_seg1_heuristic/extrap.tex
\begin{frame}
\onslide<1->{\frametitle{Extrapolation (Prediction)}}
 \begin{itemize}
\onslide<2->{ \item In the running average example, the dynamics were trivial,
 i.e., $x(k+1)=x(k)$}
\onslide<3->{ \item This means that the state never really changes.  It's
 always equal to its previous value.}
\onslide<4->{ \item We just don't know what the exact value is.}
\onslide<5->{ \item Example:  Measuring the same resistor over and over again with a voltmeter to
 estimate the resistor's resistance value.  Or getting on and off a scale repeatedly.}
\end{itemize}
\end{frame}

\begin{frame}
\onslide<1->{\frametitle{Extrapolation (Prediction)}}
 \begin{itemize}
 \onslide<2->{\item In more general Kalman filtering, we complicate the
 situation by adding two more elements:}
 \begin{itemize}
\onslide<3->{ \item State dynamics with time dependence and process noise: 
\begin{center}$x(k+1)=f\left[k,x(k),w(k)\right]$\end{center}}
\onslide<4->{ \item Typically, the process noise is additive: 
\begin{center}$x(k+1)=f\left[k,x(k)\right]+w(k)$\end{center}}
 \end{itemize}
\end{itemize}
\end{frame}

\begin{frame}
\onslide<1->{\frametitle{Extrapolation (Prediction)}}
\begin{itemize}
  \onslide<2->{\item Form of function $f$ can be simple or complicated, linear or nonlinear}
  \onslide<3->{\item Independent variable $k$ indicates  possible explicit time dependence}
  \onslide<4->{\item Process noise assumed to be Gaussian white noise (to be defined later)}
\end{itemize}
\end{frame}

\begin{frame}
\onslide<1->{\frametitle{Extrapolation (Prediction)}}
 \begin{itemize}
 \onslide<2->{\item General Kalman filtering further complicates things by
 allowing the state variable to be vector-valued}
 \onslide<3->{\item This turns the function $f$ into a vector-valued function,
 often a matrix (possibly time-dependent)}
 \onslide<4->{\item Process noise must therefore also become a vector-valued variable}
 \end{itemize}
\end{frame}

\begin{frame}
\onslide<1->{\frametitle{Extrapolation (Prediction)}}
\begin{itemize}
	\onslide<2->{\item Note that presence of a random forcing function like
	 Gaussian white noise makes the state vector process a stochastic process (aka random process)}
	  \begin{itemize} 
  		\onslide<3->{ \item If white noise were zero, problem would be deterministic}
	  \end{itemize}
	 \onslide<4->{\item We'll see what this means when we take a look at random variables and random processes}
	 \onslide<5->{\item Even for the same initial conditions the sample trajectory will be different each time, due to 
	  randomly changing forcing functions}
\end{itemize}
\end{frame}

\begin{frame}
\onslide<1->{\frametitle{Extrapolation (Prediction)}}
 \begin{itemize}
  \onslide<2->{\item Note that we did not mention a control vector $u$ that could also be part of the 
forcing function}
\begin{itemize}
\onslide<3->{\item Control vectors usually considered to be deterministic although in reality you can't control things 
perfectly! }
\onslide<4->{\item E.g., aircraft actuator (e.g., elevator or flaps) deflection angle might be a 
control quantity yet its precise value may not be known!)}
\end{itemize}
 \onslide<5->{\item We are going to ignore the presence of deterministic control and focus solely on state estimation}
\end{itemize}
\end{frame}
